\subsection{散度算子及相关空间}
\label{sec:chap4-divergence-spaces}

\renewcommand{\thestatement}{3.\arabic{statement}}
\setcounter{statement}{0}
\newcounter{chapterfoursectionthreeremark}
\renewcommand{\thechapterfoursectionthreeremark}{3.\arabic{chapterfoursectionthreeremark}}

\subsubsection{散度的右逆}
\label{subsec:chap4-divergence-right-inverse}

上述结果蕴含: 从 $(H_0^1(\Omega))^d$ 到 $L_0^2(\Omega)$ 定义的散度算子具有连续右逆.

\begin{Theorem}
\label{thm:chap4-divergence-right-inverse}
设 $\Omega$ 是 $\mathbb R^d$ 中连通、有界的 Lipschitz 区域. 存在从 $L_0^2(\Omega)$ 到 $(H_0^1(\Omega))^d$ 的连续线性算子 $\Pi$, 使得对所有 $q\in L_0^2(\Omega)$, 函数 $u=\Pi(q)$ 满足
\[
 \operatorname{div}u=q.
\]
\end{Theorem}

\par\noindent\refstepcounter{chapterfoursectionthreeremark}\textbf{Remark~\thechapterfoursectionthreeremark:}\label{rem:chap4-bogovskii-right-inverse}\quad 注意, 利用奇异积分理论, 这个结果还有一个略为显式的证明. 这是由 Bogovskii 给出的构造 (见\MBUnresolvedReference{bib:bogovskii-right-inverse}{[18, 19, 63, 93]}).

下面还会给出光滑二维区域情形中的一个较简单证明.

\par\noindent\refstepcounter{chapterfoursectionthreeremark}\textbf{Remark~\thechapterfoursectionthreeremark:}\label{rem:chap4-divergence-right-inverse-lp}\quad 按照 Remark~\ref{rem:chap4-necas-lq-version}, 可以证明: 若把 $L_0^2(\Omega)$ 换成 $L_0^p(\Omega)$, 并把 $(H_0^1(\Omega))^d$ 换成 $(W_0^{1,p}(\Omega))^d$, 则对 $1<p<+\infty$ 同一结果仍然成立. 然而, 对 $p=1$ 或 $p=+\infty$ 该结果不成立. 关于这些问题的详细讨论, 特别是上述定理的某种替代证明, 例如可参见\MBUnresolvedReference{bib:divergence-right-inverse-lp}{[21]}.

\begin{Proof}
先前已经看到 (Theorem~\ref{thm:chap4-gradient-image-closed}), 从 $L^2(\Omega)$ 到 $(H^{-1}(\Omega))^d$ 的映射 $\nabla$ 具有闭值域, 因而该值域是 Banach 空间. 令 $N$ 为这个映射的核; 则映射 $\nabla$ 可以看作从 $L^2(\Omega)/N$ 到 $\nabla(L^2(\Omega))$ 的连续线性双射. 因此由开映射定理 (Theorem~\ref{thm:chap2-open-mapping}), $\nabla$ 给出 $L^2(\Omega)/N$ 与 $\nabla(L^2(\Omega))$ 之间的 Banach 空间同构.

所以该映射的伴随给出从 $(\nabla(L^2(\Omega)))'$ 到 $(L^2(\Omega)/N)'$ 的同构. 下面辨认所涉及的对象:
\begin{itemize}
 \item $(\nabla(L^2(\Omega)))'$ 是 $\nabla(L^2(\Omega))$ 上连续线性泛函的集合; 它自然地等同于 $((H^{-1}(\Omega))^d)'$ 关于所有在 $\nabla(L^2(\Omega))$ 上为零的线性泛函所成空间的商, 使用 Theorem~\ref{thm:chap4-de-rham-hminusone} 证明中的记号, 后一个空间就是 $Z$. 此外, $(H_0^1(\Omega))^d$ 自反, 所以 $((H^{-1}(\Omega))^d)'$ 自然地等同于 $(H_0^1(\Omega))^d$. 总之,
 \[
  (\nabla(L^2(\Omega)))'=(H_0^1(\Omega))^d/Z.
 \]

 \item $(L^2(\Omega)/N)'$ 的元素是 $L^2(\Omega)$ 上在 $N$ 上为零的线性泛函; 由于 $\Omega$ 连通, $N$ 是常函数集合. $L^2(\Omega)$ 上的连续线性泛函由某个函数 $q\in L^2(\Omega)$ 表示. 它在常数上为零, 当且仅当
 \[
  \int_\Omega q(x)\,\mathrm dx=0.
 \]
 因此
 \[
  (L^2(\Omega)/N)'
  =\left\{q\in L^2(\Omega),\ \int_\Omega q(x)\,\mathrm dx=0\right\}
  =L_0^2(\Omega).
 \]

 \item 最后, 对所有 $p\in L^2(\Omega)$ 和 $u\in(H_0^1(\Omega))^d$, 分部积分给出
 \[
  \langle\nabla p,u\rangle_{H^{-1},H_0^1}
  =-(p,\operatorname{div}u)_{L^2(\Omega)},
 \]
 这说明算子 $\nabla$ 的伴随是算子 $(-\operatorname{div})$.
\end{itemize}

总之, 已经证明算子 $(-\operatorname{div})$ 给出 $(H_0^1(\Omega))^d/Z$ 与 $L_0^2(\Omega)$ 之间的同构. 记其逆同构为 $\Pi_0$. 由于 $Z$ 是 $(H_0^1(\Omega))^d$ 的闭线性子空间, 已知到 $Z^\perp$ 上的正交投影算子 $P_{Z^\perp}$ 给出从 $(H_0^1(\Omega))^d/Z$ 到 $Z^\perp$ 的同构. 因此, 从 $L_0^2(\Omega)$ 到 $(H_0^1(\Omega))^d$ 的算子
\[
 \Pi=P_{Z^\perp}\circ\Pi_0
\]
满足所需性质.
\end{Proof}

上述结果蕴含如下著名的 inf-sup 不等式 (也称为 Ladyzhenskaya--Babuška--Brezzi (LBB) 不等式):
\[
 \inf_{\substack{p\in L_0^2(\Omega)\\p\neq0}}
 \left(
  \sup_{\substack{v\in(H_0^1(\Omega))^d\\v\neq0}}
  \frac{\displaystyle\int_\Omega p(\operatorname{div}v)\,\mathrm dx}
       {\|p\|_{L^2}\|v\|_{H^1}}
 \right)>0.
\]
事实上, 该不等式等价于从 $L_0^2(\Omega)$ 到 $(H_0^1(\Omega))^d$ 的散度算子存在连续右逆.

相关不等式在 Stokes 问题的数值分析中特别重要 (例如见\MBUnresolvedReference{bib:stokes-inf-sup-numerics}{[65, 55]}).

现在给出 Theorem~\ref{thm:chap4-divergence-right-inverse} 在一个特殊情形下的替代证明, 例如见\MBUnresolvedReference{bib:divergence-right-inverse-2d}{[26]}.

\par\vspace{0.1cm}\noindent\textbf{Proof (in the case where $\Omega$ is a domain of $\mathbb R^2$ of class $C^{1,1}$):}\quad
在 $\mathbb R^2$ 上选取通常的逆时针定向, 并对任意两个向量 $a,b\in\mathbb R^2$ 定义
\[
 a\wedge b=a_1b_2-a_2b_1\in\mathbb R.
\]
设 $q\in L_0^2(\Omega)$, 并考虑如下 Laplace--Neumann 问题:
\begin{equation}
 \begin{cases}
  -\Delta\psi=q,&\text{in }\Omega,\\
  -\dfrac{\partial\psi}{\partial\nu}=0,&\text{on }\partial\Omega.
 \end{cases}
 \label{eq:chap4-laplace-neumann-right-inverse}
\end{equation}
由于 $\Omega$ 属于 $C^{1,1}$ 类且 $\int_\Omega q\,\mathrm dx=0$, Theorem~\ref{thm:chap3-laplace-neumann} 蕴含式\eqref{eq:chap4-laplace-neumann-right-inverse}存在解 $\psi\in H^2(\Omega)$, 并且 $\|\psi\|_{H^2}\leq C\|q\|_{L^2}$. 令 $v=-\nabla\psi$, 则 $v\in(H^1(\Omega))^2$, $\operatorname{div}v=q$, 且 $\|v\|_{H^1}\leq C\|q\|_{L^2}$. 遗憾的是, $v$ 在边界上并不为零, 因为只有
\[
 v\cdot\nu=-\frac{\partial\psi}{\partial\nu}=0.
\]

在 $\partial\Omega$ 上令 $g=v\wedge\nu$, 则 $g\in H^{1/2}(\partial\Omega)$. 现在使用 Remark~\ref{rem:chap3-normal-lifting-decomposition} 中构造的提升算子 $R$. 令 $\varphi=R(0,g)\in H^2(\Omega)$, 则 $\varphi$ 满足
\[
 \|\varphi\|_{H^2}
 \leq C\|g\|_{H^{1/2}(\partial\Omega)}
 \leq C'\|v\|_{H^1}
 \leq C''\|q\|_{L^2},
\]
\[
 \gamma_0(\varphi)=0,
 \qquad
 \gamma_\nu(\varphi)=\frac{\partial\varphi}{\partial\nu}=g.
\]
现在引入
\[
 w=(\nabla\varphi)^\perp
 =\left(-\frac{\partial\varphi}{\partial x_2},
 \frac{\partial\varphi}{\partial x_1}\right)
 \in(H^1(\Omega))^2.
\]
\begin{itemize}
 \item 由于 $\varphi$ 在 $\partial\Omega$ 上为常数, 有 $\nabla_T\varphi=0$ 于 $\partial\Omega$. 因此
 \[
  \nabla\varphi=(\nu\cdot\nabla\varphi)\nu,
  \qquad\text{on }\partial\Omega,
 \]
 所以 $w=(\nabla\varphi)^\perp$ 在边界上与 $\nu$ 正交; 即 $w\cdot\nu=0$ 于 $\partial\Omega$.
 \item 按构造,
 \[
  w\wedge\nu=(\nabla\varphi)^\perp\wedge\nu
  =-(\nabla\varphi)\cdot\nu=-g,
 \]
 且 $\operatorname{div}w=0$.
\end{itemize}
汇集 $v$ 与 $w$ 的性质, 可见 $u=v+w$ 满足
\[
 \operatorname{div}u=\operatorname{div}(v+w)=q,
\]
\[
 \|u\|_{H^1}\leq\|v\|_{H^1}+\|w\|_{H^1}
 \leq C\|q\|_{L^2},
\]
\[
 u\cdot\nu=v\cdot\nu+w\cdot\nu=0,
 \qquad
 u\wedge\nu=v\wedge\nu+w\wedge\nu=0.
\]
所以 $u\in(H_0^1(\Omega))^2$; 由于 $u$ 的定义线性地依赖于 $q$, 结论得证.
\par\hfill$\blacksquare$\par\vspace{0.1cm}

\subsubsection{空间 $H_{\operatorname{div}}(\Omega)$}
\label{subsec:chap4-hdiv-space}

\begin{Definition}
\label{def:chap4-hdiv-space}
设 $\Omega$ 是 $\mathbb R^d$ 中边界紧的 Lipschitz 区域. 引入空间
\[
 H_{\operatorname{div}}(\Omega)
 =\{u\in(L^2(\Omega))^d,\ \operatorname{div}u\in L^2(\Omega)\},
\]
并赋予图范数
\[
 u\longmapsto
 \bigl(\|u\|_{L^2}^2+\|\operatorname{div}u\|_{L^2}^2\bigr)^{1/2}.
\]
此外, 定义 $H_{\operatorname{div},0}(\Omega)$ 为 $(\mathcal D(\Omega))^d$ 在 $H_{\operatorname{div}}(\Omega)$ 中的闭包.
\end{Definition}

使用 Chapter~\ref{chap:sobolev-spaces} 的 Section~\ref{subsec:chap3-first-order-operator-domains} 中引入的记号, 可见 $H_{\operatorname{div}}(\Omega)$ 恰好是散度算子在 $L^2$ 中的定义域. 更准确地说,
\[
 H_{\operatorname{div}}(\Omega)=W_{\operatorname{div}}^{2,2}(\Omega).
\]

于是可以用这个记号重新表述 Chapter~\ref{chap:sobolev-spaces} 的 Section~\ref{subsec:chap3-first-order-operator-domains} 中得到的各个结果:
\begin{itemize}
 \item 空间 $(C_c^\infty(\overline\Omega))^d$ 在 $H_{\operatorname{div}}(\Omega)$ 中稠密 (见 Theorem~\ref{thm:chap3-graph-space-density}).
 \item 存在从 $H_{\operatorname{div}}(\Omega)$ 到 $H^{-1/2}(\partial\Omega)$ 的连续迹算子 $\gamma_\nu$, 使对任意 $u\in(C_c^\infty(\overline\Omega))^d$ 有 $\gamma_\nu(u)=u\cdot\nu$. 此外, 成立如下 Stokes 公式:
 \begin{equation}
  \int_\Omega u\cdot\nabla w\,\mathrm dx
  +\int_\Omega w\operatorname{div}u\,\mathrm dx
  =\langle\gamma_\nu(u),\gamma_0(w)\rangle_{H^{-1/2},H^{1/2}},
  \quad\forall u\in H_{\operatorname{div}}(\Omega),\ \forall w\in H^1(\Omega).
  \label{eq:chap4-hdiv-stokes-formula}
 \end{equation}
 这来自 Theorem~\ref{thm:chap3-generalized-normal-trace}.
 \item 空间 $H_{\operatorname{div},0}(\Omega)$ 是迹算子 $\gamma_\nu$ 的核 (见 Theorem~\ref{thm:chap3-graph-space-zero-trace}).
\end{itemize}

最后陈述上述结果到张量场的推广. 把 $(H_{\operatorname{div}}(\Omega))^d$ 与满足散度属于 $(L^2(\Omega))^d$ 的张量 $\sigma\in(L^2(\Omega))^{d\times d}$ 等同. 逐行推理, 可以定义仍记为
\[
 \gamma_\nu:(H_{\operatorname{div}}(\Omega))^d
 \longrightarrow(H^{-1/2}(\partial\Omega))^d
\]
的连续线性算子, 使对任意光滑张量场 $\sigma$ 有 $\gamma_\nu(\sigma)=\sigma\cdot\nu$. 立即得到如下结果.

\begin{Lemma}
\label{lem:chap4-tensor-stokes-formula}
设 $\Omega$ 是 $\mathbb R^d$ 中有界的 Lipschitz 区域. 对任意 $v\in(H^1(\Omega))^d$ 和任意 $\sigma\in(H_{\operatorname{div}}(\Omega))^d$, 成立如下 Stokes 公式:
\[
 \int_\Omega\sigma:\nabla v\,\mathrm dx
 +\int_\Omega v\cdot\operatorname{div}\sigma\,\mathrm dx
 =\langle\sigma\cdot\nu,v\rangle_{H^{-1/2},H^{1/2}}.
\]
\end{Lemma}

\subsubsection{无散度向量场. Leray 分解}
\label{subsec:chap4-leray-decomposition}

考虑紧支撑于 $\Omega$ 的光滑无散度向量场集合 $\mathcal V$:
\[
 \mathcal V=\{\varphi\in(\mathcal D(\Omega))^d,\ \operatorname{div}\varphi=0\}.
\]
于是可以定义如下基本空间: $V$ 为 $\mathcal V$ 在 $(H_0^1(\Omega))^d$ 中的闭包, $H$ 为 $\mathcal V$ 在 $(L^2(\Omega))^d$ 中的闭包. 这两个空间当然都是 Hilbert 空间, 分别赋予由 $(H_0^1(\Omega))^d$ 和 $(L^2(\Omega))^d$ 的内积诱导的内积.

现在可以更精确地刻画这两个空间.

\begin{Lemma}
\label{lem:chap4-divergence-free-h01-closure}
设 $\Omega$ 是 $\mathbb R^d$ 中有界的 Lipschitz 区域. 则
\[
 V=\{v\in(H_0^1(\Omega))^d,\ \operatorname{div}v=0\}.
\]
\end{Lemma}

\begin{Proof}
把空间 $\{v\in(H_0^1(\Omega))^d,\ \operatorname{div}v=0\}$ 记为 $\widetilde V$. 需要证明 $V=\widetilde V$.

显然 $\mathcal V\subset\widetilde V$, 且 $\widetilde V$ 在 $(H_0^1(\Omega))^d$ 中闭, 所以包含关系 $V\subset\widetilde V$ 显然成立. 为证明另一个包含关系, 使用对偶论: 证明 $\widetilde V$ 上任意在 $\mathcal V$ 上为零的连续线性泛函也在 $\widetilde V$ 上为零 (根据 Proposition~\ref{prop:chap2-density-by-annihilator}, 这已经充分).

因此, 若 $f\in(H^{-1}(\Omega))^d\cap\mathcal V^\perp$, 则由 Theorem~\ref{thm:chap4-de-rham}, $f$ 是某个函数 $p\in L_0^2(\Omega)$ 的梯度. 于是对任意 $v\in\widetilde V$,
\[
 \langle f,v\rangle_{H^{-1},H_0^1}
 =\langle\nabla p,v\rangle_{H^{-1},H_0^1}
 =-(p,\operatorname{div}v)_{L^2}=0.
\]
所以已经证明 $f$ 在 $\widetilde V$ 上为零, 证明完成.
\end{Proof}

\begin{Theorem}
\label{thm:chap4-leray-decomposition-characterization}
设 $\Omega$ 是 $\mathbb R^d$ 中连通、有界的 Lipschitz 区域. 则
\[
 H=\{u\in(L^2(\Omega))^d,\ \operatorname{div}u=0,\ \gamma_\nu(u)=0\},
\]
并且 $H$ 在 $(L^2(\Omega))^d$ 中的正交补满足
\begin{equation}
 H^\perp=\{u\in(L^2(\Omega))^d,\ \exists p\in H^1(\Omega),\ u=\nabla p\}.
 \label{eq:chap4-leray-orthogonal-gradient-space}
\end{equation}
\end{Theorem}

\begin{Proof}
\begin{itemize}
 \item 首先证明刻画 $H$ 在 $(L^2(\Omega))^d$ 中正交补的第二个等式. 记
 \[
  \mathcal H^*=\{u\in(L^2(\Omega))^d,\ \exists p\in H^1(\Omega),\ u=\nabla p\}.
 \]
 设 $u\in\mathcal H^*$; 则对所有 $v\in\mathcal V$,
 \[
  (u,v)_{L^2}
  =\int_\Omega u\cdot v\,\mathrm dx
  =\int_\Omega\nabla p\cdot v\,\mathrm dx
  =-\int_\Omega p\operatorname{div}v\,\mathrm dx=0.
 \]
 由 $\mathcal V$ 在 $H$ 中关于 $(L^2(\Omega))^d$ 拓扑的稠密性, 这说明 $u\in H^\perp$, 因而 $\mathcal H^*\subset H^\perp$. 反过来, 若 $u\in H^\perp$, 则特别地, 对所有 $v\in\mathcal V$,
 \[
  \langle u,v\rangle_{H^{-1},H_0^1}=(u,v)_{L^2}=0.
 \]
 所以由 Theorem~\ref{thm:chap4-de-rham}, 存在 $p\in L_0^2(\Omega)$, 使得 $u=\nabla p$. 由于 $u\in(L^2(\Omega))^d$, 的确有 $p\in H^1(\Omega)$. 因而 $u\in\mathcal H^*$.

 \item 现在证明第一个等式. 记
 \[
  \mathcal H=\{u\in(L^2(\Omega))^d,\ \operatorname{div}u=0,\ \gamma_\nu(u)=0\}.
 \]
 先证明包含关系 $H\subset\mathcal H$. 若 $u\in H$, 则存在 $\mathcal V$ 中的序列 $(u_n)_n$, 它在 $(L^2(\Omega))^d$ 的拓扑下收敛到 $u$. 函数 $u_n$ 无散度, 因而它的极限 $u$ 也无散度 (在 $\mathcal D'(\Omega)$ 中). 所以 $u\in H_{\operatorname{div}}(\Omega)$, 且
 \[
  \lim_{n\to\infty}\|u-u_n\|_{H_{\operatorname{div}}(\Omega)}=0.
 \]
 由法向迹算子 $\gamma_\nu$ 的连续性, 推出
 \[
  \lim_{n\to\infty}
  \|\gamma_\nu(u)-\gamma_\nu(u_n)\|_{H^{-1/2}(\partial\Omega)}=0,
 \]
 所以 $\gamma_\nu(u)=0$. 这恰好说明 $u\in\mathcal H$.

 现在已知 $H\subset\mathcal H$, 因而把 $H$ 在 $\mathcal H$ 中的正交补记为 $\mathcal H^{**}$. 设 $u\in\mathcal H^{**}\subset H^\perp$. 由式\eqref{eq:chap4-leray-orthogonal-gradient-space}的刻画, 存在 $p\in H^1(\Omega)$ 使 $u=\nabla p$; 又由于 $u\in\mathcal H$, 可见 $p$ 满足问题
 \[
  \begin{cases}
   \Delta p=\operatorname{div}u=0,&\text{in }\Omega,\\
   \dfrac{\partial p}{\partial\nu}=\gamma_\nu(u)=0,&\text{on }\partial\Omega.
  \end{cases}
 \]
 使用式\eqref{eq:chap4-hdiv-stokes-formula}, 得到
 \[
  \int_\Omega|\nabla p|^2\,\mathrm dx=0,
 \]
 因而 $u=\nabla p=0$. 所以 $\mathcal H^{**}=\{0\}$. 因此 $\mathcal H=H$.
\end{itemize}
\end{Proof}

因此, 空间 $(L^2(\Omega))^d$ 正交分解成两个子空间: 散度为零且法向迹 $\gamma_\nu(u)$ 为零的向量场空间 $H$, 以及 $H^1(\Omega)$ 中函数的梯度场空间 $H^\perp$.

\begin{Definition}
\label{def:chap4-leray-decomposition}
分解
\[
 (L^2(\Omega))^d=H\mathbin{\mathop\oplus\limits^\perp}H^\perp
\]
称为空间 $(L^2(\Omega))^d$ 的 Leray 分解. 从 $(L^2(\Omega))^d$ 到 $H$ 的正交投影称为 Leray 投影, 记为 $\mathcal P$.
\end{Definition}

当在全空间 $\Omega=\mathbb R^d$ 中工作时, 可以在 Fourier 变量下给出 Leray 投影的显式表达式. 遗憾的是, 在有界开集中, 这种分析要复杂得多. 特别地, 由于边界条件, 容易看出, 与全空间或周期情形不同, Leray 投影 $\mathcal P$ 不与通常的微分算子交换. 此外, 算子 $\mathcal P$ 是非局部的; $\mathcal Pu$ 的值依赖于函数 $u$ 在整个 $\Omega$ 上的值, 而不仅依赖于所考察点邻域中的值.

当区域足够光滑时, 可以证明 Leray 投影 $\mathcal P$ 把 $H^1(\Omega)^d$ 映入自身. 更准确地说, 有如下结果.

\begin{Proposition}
\label{prop:chap4-leray-projection-h1}
假设 $\Omega$ 是 $C^{1,1}$ 类有界区域. 则对任意 $u\in(H^1(\Omega))^d$, 有 $\mathcal Pu\in(H^1(\Omega))^d$, 且
\[
 \|\mathcal Pu\|_{H^1}\leq C\|u\|_{H^1},
\]
其中 $C$ 仅依赖于 $\Omega$.
\end{Proposition}

\begin{Proof}
按定义, 写成 $u=\mathcal Pu+\nabla p$, 其中 $p\in H^1(\Omega)$, $\operatorname{div}(\mathcal Pu)=0$, 且 $\gamma_\nu(\mathcal Pu)=0$. 对这个等式取散度, 得到
\[
 \Delta p=\operatorname{div}u\in L^2(\Omega),
\]
因为假设了 $u\in(H^1(\Omega))^d$. 此外,
\[
 \gamma_\nu(\nabla p)
 =\gamma_\nu(u)+\gamma_\nu(\mathcal Pu)
 =\gamma_\nu(u)=\gamma_0(u)\cdot\nu
 \in H^{1/2}(\partial\Omega).
\]
因此, $p$ 是具有 $L^2(\Omega)$ 中源项和 $H^{1/2}(\partial\Omega)$ 中 Neumann 边界数据的 Laplace 方程之解. Laplace--Neumann 问题的椭圆正则性 (见 Theorem~\ref{thm:chap3-laplace-neumann}) 蕴含 $p\in H^2(\Omega)$, 并且
\[
 \|p\|_{H^2}
 \leq C\bigl(\|\operatorname{div}u\|_{L^2}
 +\|\gamma_0(u)\cdot\nu\|_{H^{1/2}}\bigr)
 \leq C'\|u\|_{H^1}.
\]
最后得到 $\mathcal Pu=u-\nabla p\in(H^1(\Omega))^d$, 且 $\|\mathcal Pu\|_{H^1}\leq C\|u\|_{H^1}$.
\end{Proof}
